\subsection{情景结果与机制讨论}
在2022公报归一化口径下，基线情景的江豚和长江鲟末值为1.133和1.140。若只增强通道阻隔、暂不叠加额外污染，二者降至0.901和0.856。把长江鲟放流项置为$R_A=0$后，长江鲟降至0.519，而江豚仅变为1.145；污染情景下，两者又降至0.610和0.704。这里改变的是通道损失参数和放流控制量；放流项位于长江鲟方程，通道损失同时进入两类物种方程。通道约束同时压低两类物种，放流变化主要落在长江鲟上。表中报告江豚和长江鲟终态，$M$与$V$留作后文解释过程差异。
\begin{table}[htbp]
\centering
\caption{问题二的结果口径、相对变化与对照说明}
\label{tab:q2_result}
\begin{tabularx}{0.95\textwidth}{>{\raggedright\arraybackslash}p{2.2cm} >{\raggedright\arraybackslash}p{3.6cm} >{\raggedright\arraybackslash}p{2.8cm} >{\raggedright\arraybackslash}p{1.7cm} X}
\toprule
项目 & 数值/情景输出 & 相对基线变化 & 是否外部验证 & 解释用途 \\
\midrule
2022 公报基准 & 江豚 1249 头、长江鲟 438 尾 & -- & 是 & 作为问题二的初值与口径基准 \\
基线情景末值 & 江豚 1.133，长江鲟 1.140 & -- & 否 & 作为禁渔与食源修复下的基线投影 \\
$R_A=0$ 对照 & 江豚 1.145，长江鲟 0.519 & 江豚 +1.0\%，长江鲟 -54.5\% & 否 & 拆分自然恢复与人工放流补偿效应 \\
高阻隔对照 & 江豚 0.901，长江鲟 0.856 & 江豚 -20.5\%，长江鲟 -24.9\% & 否 & 仅提高通道阻隔，分离生境约束效应 \\
污染情景末值 & 江豚 0.610，长江鲟 0.704 & 江豚 -46.2\%，长江鲟 -38.2\% & 否 & 比较污染叠加阻隔对基线收益的侵蚀 \\
\bottomrule
\end{tabularx}
\end{table}
这些对照对应的机制各有差异。无放流时，长江鲟大幅下降而江豚变化很小，说明人工补偿主要作用于鲟类而不是江豚；高阻隔与污染情景同时压低两类物种，说明连通性和环境质量属于共性约束。作为定性验证，模型在$t=3$（对应2024年）时已有$A(3)/A(0)>1$，与题目所述“2024 年长江鲟自然繁殖群体重新出现产卵场”这一事实方向一致。

